%! Author = sriadityavedantam
%! Date = 3/30/26

\section{Chapter 4}

\begin{problem}{
    Let $d$ be an integer that is not a perfect square.
    Show that
    \[
        \mathbb{Q}\left(\sqrt{d}\right)=\{a+b\sqrt{d}:a,b\in\mathbb{Q}\}
    \]
    is a subfield of $\mathbb{C}$.
}
    A field is a commutative ring with identity in which every nonzero element has a multiplicative inverse.
    Thus we must show that this set is closed under addition, additive inverses, multiplication, and multiplicative inverses for nonzero elements.
    Let
    \[
        q=a_1+b_1\sqrt{d}
        \qquad\text{and}\qquad
        p=a_2+b_2\sqrt{d}.
    \]
    Then
    \begin{align*}
        q-p
        &=
        (a_1+b_1\sqrt{d})-(a_2+b_2\sqrt{d})\\
        &=
        (a_1-a_2)+(b_1-b_2)\sqrt{d}.
    \end{align*}
    Since $a_1-a_2\in\mathbb{Q}$ and $b_1-b_2\in\mathbb{Q}$, this shows closure under subtraction, and hence also under addition and additive inverses.
    Now check multiplication:
    \begin{align*}
        qp
        &=
        (a_1+b_1\sqrt{d})(a_2+b_2\sqrt{d})\\
        &=
        a_{1}a_2+a_{1}b_2\sqrt{d}+a_{2}b_1\sqrt{d}+b_{1}b_{2}d\\
        &=
        (a_{1}a_2+b_{1}b_{2}d)+(a_{1}b_2+a_{2}b_1)\sqrt{d}.
    \end{align*}
    Since $d\in\mathbb{Z}\subseteq\mathbb{Q}$, the coefficients are rational.
    Thus $\mathbb{Q}\left(\sqrt{d}\right)$ is closed under multiplication.
    Also,
    \[
        0=0+0\sqrt{d}\in\mathbb{Q}\left(\sqrt{d}\right)
    \]
    and
    \[
        1=1+0\sqrt{d}\in\mathbb{Q}\left(\sqrt{d}\right).
    \]
    Now let
    \[
        w=a+b\sqrt{d}\in\mathbb{Q}\left(\sqrt{d}\right),
        \qquad
        w\neq 0.
    \]
    We must show that $w^{-1}\in\mathbb{Q}\left(\sqrt{d}\right)$.
    Compute:
    \begin{align*}
        \frac{1}{a+b\sqrt{d}}
        &=
        \frac{1}{a+b\sqrt{d}}\cdot \frac{a-b\sqrt{d}}{a-b\sqrt{d}}\\
        &=
        \frac{a-b\sqrt{d}}{a^2-b^{2}d}\\
        &=
        \frac{a}{a^2-b^{2}d}+\frac{-b}{a^2-b^{2}d}\sqrt{d}.
    \end{align*}
    So it remains to show that
    \[
        a^2-b^{2}d\neq 0.
    \]
    If
    \[
        a^2-b^{2}d=0,
    \]
    then
    \[
        a^2=b^{2}d.
    \]
    If $b=0$, then $a=0$, which would give $w=0$, a contradiction.
    So $b\neq 0$.
    Then
    \[
        d=\left(\frac{a}{b}\right)^2.
    \]
    This would mean $d$ is a square in $\mathbb{Q}$.
    Since $d\in\mathbb{Z}$, that would imply $d$ is a perfect square in $\mathbb{Z}$, contradicting the hypothesis.
    Thus
    \[
        a^2-b^{2}d\neq 0,
    \]
    and so $w^{-1}\in\mathbb{Q}\left(\sqrt{d}\right)$.
    Therefore $\mathbb{Q}\left(\sqrt{d}\right)$ is a subfield of $\mathbb{C}$.
\end{problem}

\begin{problem}{
    If $u$ is a unit in a commutative ring $R$ with identity, prove that $u$ is not a zero divisor.
}
    Suppose $u$ is a unit.
    Then there exists $u^{-1}\in R$ such that
    \[
        uu^{-1}=1.
    \]
    Now suppose, for contradiction, that $u$ is also a zero divisor.
    Then there exists $x\in R$, $x\neq 0$, such that
    \[
        ux=0.
    \]
    Multiply both sides by $u^{-1}$:
    \begin{align*}
        u^{-1}(ux) &= u^{-1}0\\
        (u^{-1}u)x &= 0\\
        1x &= 0\\
        x &= 0.
    \end{align*}
    This contradicts the fact that $x\neq 0$.
    Therefore $u$ is not a zero divisor.
\end{problem}

\begin{problem}{
    Which of the following functions are homomorphisms?

    (a)
    \[
        f:\mathbb{Z}\mapsto\mathbb{Z},
        \qquad
        f(x)=-x.
    \]

    (e)
    \[
        f:\mathbb{Z}_{12}\mapsto\mathbb{Z}_4
    \]
    defined by
    \[
        f([x]_{12})=[x]_4.
    \]
}
    \textbf{(a).}
    We check whether this function preserves addition and multiplication.
    For addition,
    \[
        f(a+b)=-(a+b)=-a-b=f(a)+f(b).
    \]
    So it preserves addition.
    For multiplication,
    \[
        f(ab)=-(ab),
    \]
    but
    \[
        f(a)f(b)=(-a)(-b)=ab.
    \]
    In general,
    \[
        -(ab)\neq ab,
    \]
    so $f$ does not preserve multiplication.
    Therefore this function is not a ring homomorphism.

    \textbf{(e).}
    We first show that the function is well-defined.
    Suppose
    \[
        [x]_{12}=[y]_{12}.
    \]
    Then
    \[
        x\equiv y\mod{12},
    \]
    so
    \[
        12\mid (x-y).
    \]
    Hence
    \[
        4\mid (x-y),
    \]
    which means
    \[
        x\equiv y\mod 4.
    \]
    Thus
    \[
        [x]_4=[y]_4.
    \]
    So the map is well-defined.
    Now check addition:
    \begin{align*}
        f([x]_{12}+[y]_{12})
        &=
        f([x+y]_{12})\\
        &=
        [x+y]_4\\
        &=
        [x]_4+[y]_4\\
        &=
        f([x]_{12})+f([y]_{12}).
    \end{align*}
    Now check multiplication:
    \begin{align*}
        f([x]_{12}[y]_{12})
        &=
        f([xy]_{12})\\
        &=
        [xy]_4\\
        &=
        [x]_4[y]_4\\
        &=
        f([x]_{12})f([y]_{12}).
    \end{align*}
    Therefore this function is a homomorphism.
\end{problem}

\begin{problem}{
    Let $R$ and $S$ be rings.

    (a) Prove that
    \[
        f:R\times S\mapsto R
    \]
    given by
    \[
        f((r,s))=r
    \]
    is a surjective homomorphism.

    (b) Prove that
    \[
        g:R\times S\mapsto S
    \]
    given by
    \[
        g((r,s))=s
    \]
    is a surjective homomorphism.

    (c) If both $R$ and $S$ are nonzero rings, prove that the homomorphisms $f$ and $g$ are not injective.
}
    \textbf{(a).}
    First we show surjectivity.
    Let $r\in R$.
    Choose any $s\in S$.
    Then
    \[
        f((r,s))=r.
    \]
    So $f$ is surjective.
    Now check addition:
    \begin{align*}
        f((r_1,s_1)+(r_2,s_2))
        &=
        f((r_1+r_2,s_1+s_2))\\
        &=
        r_1+r_2\\
        &=
        f((r_1,s_1))+f((r_2,s_2)).
    \end{align*}
    Now check multiplication:
    \begin{align*}
        f((r_1,s_1)(r_2,s_2))
        &=
        f((r_{1}r_2,s_{1}s_2))\\
        &=
        r_{1}r_2\\
        &=
        f((r_1,s_1))f((r_2,s_2)).
    \end{align*}
    Thus $f$ is a surjective homomorphism.

    \textbf{(b).}
    First we show surjectivity.
    Let $s\in S$.
    Choose any $r\in R$.
    Then
    \[
        g((r,s))=s.
    \]
    So $g$ is surjective.
    Now check addition:
    \begin{align*}
        g((r_1,s_1)+(r_2,s_2))
        &=
        g((r_1+r_2,s_1+s_2))\\
        &=
        s_1+s_2\\
        &=
        g((r_1,s_1))+g((r_2,s_2)).
    \end{align*}
    Now check multiplication:
    \begin{align*}
        g((r_1,s_1)(r_2,s_2))
        &=
        g((r_{1}r_2,s_{1}s_2))\\
        &=
        s_{1}s_2\\
        &=
        g((r_1,s_1))g((r_2,s_2)).
    \end{align*}
    Thus $g$ is a surjective homomorphism.

    \textbf{(c).}
    Since both $R$ and $S$ are nonzero rings, choose
    \[
        s_1,s_2\in S
        \qquad\text{with}\qquad
        s_1\neq s_2,
    \]
    and choose any $r\in R$.
    Then
    \[
        f((r,s_1))=r=f((r,s_2)),
    \]
    but
    \[
        (r,s_1)\neq (r,s_2).
    \]
    So $f$ is not injective.
    Similarly, choose
    \[
        r_1,r_2\in R
        \qquad\text{with}\qquad
        r_1\neq r_2,
    \]
    and choose any $s\in S$.
    Then
    \[
        g((r_1,s))=s=g((r_2,s)),
    \]
    but
    \[
        (r_1,s)\neq (r_2,s).
    \]
    So $g$ is not injective.
\end{problem}

\begin{problem}{
    Show that the first ring is not isomorphic to the second.

    (e)
    \[
        \mathbb{Z}\times\mathbb{Z}_2
        \qquad\text{and}\qquad
        \mathbb{Z}.
    \]

    (f)
    \[
        \mathbb{Z}_4\times\mathbb{Z}_4
        \qquad\text{and}\qquad
        \mathbb{Z}_{16}.
    \]
}
    \textbf{(e).}
    In
    \[
        \mathbb{Z}\times\mathbb{Z}_2,
    \]
    the element
    \[
        (0,[1])
    \]
    is a nontrivial idempotent, since
    \[
        (0,[1])^2=(0,[1]).
    \]
    But in $\mathbb{Z}$, the only idempotent elements are $0$ and $1$, because if
    \[
        x^2=x,
    \]
    then
    \[
        x(x-1)=0,
    \]
    so
    \[
        x=0
        \quad\text{or}\quad
        x=1.
    \]
    A ring isomorphism preserves idempotents.
    Since $\mathbb{Z}\times\mathbb{Z}_2$ has a nontrivial idempotent different from $(0,[0])$ and $(1,[1])$, while $\mathbb{Z}$ does not, the rings are not isomorphic.

    \textbf{(f).}
    In
    \[
        \mathbb{Z}_4\times\mathbb{Z}_4,
    \]
    the element
    \[
        ([2],[0])
    \]
    has additive order $2$, since
    \[
        ([2],[0])+([2],[0])=([0],[0]).
    \]
    So this ring has nonzero elements of additive order $2$.
    In $\mathbb{Z}_{16}$, the additive group is cyclic of order $16$, and it has exactly one element of additive order $2$, namely $[8]$.
    In contrast, $\mathbb{Z}_4\times\mathbb{Z}_4$ has three nonzero elements of additive order $2$:
    \[
        ([2],[0]),\quad ([0],[2]),\quad ([2],[2]).
    \]
    An isomorphism of rings induces an isomorphism of additive groups, so it must preserve the number of elements of each order.
    Therefore
    \[
        \mathbb{Z}_4\times\mathbb{Z}_4
        \not\cong
        \mathbb{Z}_{16}.
    \]
\end{problem}

\begin{problem}{
    Let $F$ be a field and
    \[
        f:F\mapsto R
    \]
    a homomorphism of rings.

    (a) If there is a nonzero element $c$ of $F$ such that
    \[
        f(c)=0_R,
    \]
    prove that $f$ is the zero homomorphism.

    (b) Prove that $f$ is either injective or the zero homomorphism.
}
    \textbf{(a).}
    Assume there is a nonzero element $c\in F$ such that
    \[
        f(c)=0_R.
    \]
    Since $F$ is a field and $c\neq 0$, $c$ is a unit.
    So there exists $c^{-1}\in F$ such that
    \[
        cc^{-1}=1.
    \]
    Now let $x\in F$.
    Then
    \begin{align*}
        f(x)
        &=
        f(x\cdot 1)\\
        &=
        f(xcc^{-1})\\
        &=
        f(x)f(c)f(c^{-1})\\
        &=
        f(x)\cdot 0_R\cdot f(c^{-1})\\
        &=
        0_R.
    \end{align*}
    Since this is true for every $x\in F$, $f$ is the zero homomorphism.

    \textbf{(b).}
    We prove the contrapositive.
    Suppose $f$ is not injective.
    Then there exist $a,b\in F$ with $a\neq b$ such that
    \[
        f(a)=f(b).
    \]
    Thus
    \[
        f(a-b)=0_R.
    \]
    Since $a\neq b$, we have
    \[
        a-b\neq 0.
    \]
    By part (a), it follows that $f$ is the zero homomorphism.
    Therefore, if $f$ is not the zero homomorphism, then $f$ must be injective.
    So $f$ is either injective or the zero homomorphism.
\end{problem}

\begin{problem}{
    Let $m,n\in\mathbb{Z}$ with $\gcd(m,n)=1$, and let
    \[
        f:\mathbb{Z}_{mn}\mapsto\mathbb{Z}_m\times\mathbb{Z}_n
    \]
    be the function given by
    \[
        f([a]_{mn})=([a]_m,[a]_n).
    \]

    (a) Show that the map $f$ is well-defined.

    (b) Prove that $f$ is an isomorphism.
}
    \textbf{(a).}
    Suppose
    \[
        [a]_{mn}=[b]_{mn}
    \]
    in $\mathbb{Z}_{mn}$.
    Then
    \[
        a\equiv b\mod {mn},
    \]
    so
    \[
        mn\mid (a-b).
    \]
    Hence
    \[
        m\mid (a-b)
        \qquad\text{and}\qquad
        n\mid (a-b).
    \]
    Therefore
    \[
        [a]_m=[b]_m
        \qquad\text{and}\qquad
        [a]_n=[b]_n.
    \]
    So
    \[
        f([a]_{mn})=([a]_m,[a]_n)=([b]_m,[b]_n)=f([b]_{mn}).
    \]
    Thus $f$ is well-defined.

    \textbf{(b).}
    First we show that $f$ is a homomorphism.
    For addition,
    \begin{align*}
        f([a]_{mn}+[b]_{mn})
        &=
        f([a+b]_{mn})\\
        &=
        ([a+b]_m,[a+b]_n)\\
        &=
        ([a]_m,[a]_n)+([b]_m,[b]_n)\\
        &=
        f([a]_{mn})+f([b]_{mn}).
    \end{align*}
    For multiplication,
    \begin{align*}
        f([a]_{mn}[b]_{mn})
        &=
        f([ab]_{mn})\\
        &=
        ([ab]_m,[ab]_n)\\
        &=
        ([a]_m,[a]_n)([b]_m,[b]_n)\\
        &=
        f([a]_{mn})f([b]_{mn}).
    \end{align*}
    Now show injectivity.
    Suppose
    \[
        f([a]_{mn})=f([b]_{mn}).
    \]
    Then
    \[
        ([a]_m,[a]_n)=([b]_m,[b]_n),
    \]
    so
    \[
        [a]_m=[b]_m
        \qquad\text{and}\qquad
        [a]_n=[b]_n.
    \]
    Thus
    \[
        m\mid (a-b)
        \qquad\text{and}\qquad
        n\mid (a-b).
    \]
    Since
    \[
        \gcd(m,n)=1,
    \]
    it follows that
    \[
        mn\mid (a-b).
    \]
    Therefore
    \[
        [a]_{mn}=[b]_{mn}.
    \]
    So $f$ is injective.
    Now show surjectivity.
    Let
    \[
        ([r]_m,[s]_n)\in \mathbb{Z}_m\times\mathbb{Z}_n.
    \]
    Since
    \[
        \gcd(m,n)=1,
    \]
    there exist $u,v\in\mathbb{Z}$ such that
    \[
        um+vn=1.
    \]
    Set
    \[
        c=svn+rum.
    \]
    Then modulo $m$, we have
    \[
        svn\equiv s(1-um)\equiv s\mod m
    \]
    and
    \[
        rum\equiv 0\mod m.
    \]
    So
    \[
        c\equiv s\mod m
    \]
    is not quite what we want.
    Instead use the standard choice
    \[
        c=r(vn)+s(um).
    \]
    Then modulo $m$,
    \[
        vn\equiv 1\mod m
        \qquad\text{and}\qquad
        um\equiv 0\mod m,
    \]
    so
    \[
        c\equiv r\mod m.
    \]
    Similarly, modulo $n$,
    \[
        um\equiv 1\mod n
        \qquad\text{and}\qquad
        vn\equiv 0\mod n,
    \]
    so
    \[
        c\equiv s\mod n.
    \]
    Hence
    \[
        f([c]_{mn})=([r]_m,[s]_n).
    \]
    So $f$ is surjective.
    Therefore $f$ is a bijective homomorphism, hence an isomorphism.
\end{problem}

\begin{problem}{
    If $\gcd(m,n)\neq 1$, prove that $\mathbb{Z}_{mn}$ is not isomorphic to $\mathbb{Z}_m\times\mathbb{Z}_n$.
}
    Let
    \[
        d=\gcd(m,n)>1.
    \]
    Then we can write
    \[
        m=dq
        \qquad\text{and}\qquad
        n=ds
    \]
    for some integers $q,s$.
    Consider the element
    \[
        [mq]_{mn}=[m]_{mn}\cdot [q]_{mn},
    \]
    but an easier choice is
    \[
        [dqs]_{mn}.
    \]
    Since
    \[
        mn=d^{2}qs,
    \]
    we have
    \[
        dqs\not\equiv 0\mod {mn}
    \]
    in general, but under the map
    \[
        [a]_{mn}\mapsto ([a]_m,[a]_n),
    \]
    we get
    \[
        [dqs]_m=[0]_m
    \]
    because
    \[
        dqs=\frac{mn}{d}
    \]
    is a multiple of $m=dq$.
    Similarly,
    \[
        [dqs]_n=[0]_n
    \]
    because it is also a multiple of $n=ds$.
    So the natural map sends a nonzero class in $\mathbb{Z}_{mn}$ to
    \[
        ([0]_m,[0]_n).
    \]
    Thus the natural map is not injective.
    Now if
    \[
        \mathbb{Z}_{mn}\cong \mathbb{Z}_m\times\mathbb{Z}_n,
    \]
    then their additive groups would also be isomorphic.
    But in $\mathbb{Z}_{mn}$, the additive group is cyclic.
    On the other hand, when $\gcd(m,n)\neq 1$, the additive group
    \[
        \mathbb{Z}_m\times\mathbb{Z}_n
    \]
    is not cyclic.
    Therefore the rings cannot be isomorphic.
    Hence, if
    \[
        \gcd(m,n)\neq 1,
    \]
    then
    \[
        \mathbb{Z}_{mn}\not\cong \mathbb{Z}_m\times\mathbb{Z}_n.
    \]
\end{problem}